\scanpage{141}
\mspage{V.13}
% archive sheet number: 140
% manuscript lines 1-2
en d'autres termes que l'on a
\[
\Phi^i_{\unsure{\nu|v}}=0\qquad\text{pour }i\leq m-d-1.
\tag*{(3.1.3)}
\]
% Everything below is crossed by long oblique cancellation strokes.
% manuscript lines 3-8
\deleted{Remarque 3.2. Supposons par exemple que $X_s$ \deleted{et $X_\eta$} soit \unsure{localement} \deleted{une} intersection complète, \deleted{\unsure{pour que $X$ soit}} \inserted{ce qui vient au même que $X$ soit loc. une intersection complète,} \deleted{\unsure{implique}} ce qui implique la même chose pour $X_\eta$. Alors (\deleted{\unsure{cf.}} SGA 1 \uline{XIV} 5.6) on peut prendre $m\geq n$, et (3.1.3) prend la forme}
% slanted; the margin note is itself crossed by a cancellation stroke
\margin{\deleted{Mais dans ce cas on peut améliorer la conclusion, en remplaçant $i\leq\unsure{n|m}-d-1$ par $i\leq n-d$.}}
\[
\deleted{\Phi^i_\nu=0\qquad\text{pour }i\leq n-d-1,}
\tag*{\deleted{(3.2.1)}}
\]
% manuscript lines 9-14
\deleted{ce qui, compte tenu de 2.7, prend la forme}
\[
\deleted{\Phi^i_\nu=0\qquad\text{si }i\notin[n-d,n+d+1]}
\tag*{\deleted{(\unsure{3.1.5|3.2.5})}}
\]
\deleted{Si $d=0$, \deleted{\unsure{par}} (cas d'un \deleted{\unsure{singularités}} lieu de non lissité \uline{fini} sur $X_0$) cela signifie}
\[
\deleted{\Phi^i_\nu=0\qquad\text{si }i\ne n,n+1.}
\tag*{\deleted{(3.1.6)}}
\]
% manuscript lines 15-24
\deleted{Du point de vue action du groupe d'inertie $I$, (3.1.3) \deleted{\unsure{signifie}} implique que cette action est triviale \deleted{sur} $H^i(X_{\unsure{\bar\eta|\bar y}},\Lambda_\nu)$ si $i<m-d-1$, donc compte tenu de 2.7) si $i\notin[m-d-1,n+d]$, donc si $m=n$ (p. ex. $X_s$ loc. une intersection complète) si $i\notin[n-d-1,n+d]$, et si de plus $d=0$ si \deleted{\unsure{$i\notin[n-1,n]$}} $i\ne n-1,n$. On comparera avec la conclusion de 2.6, plus fort (car il suffit de supposer $i\ne n$) mais avec une hypothèse \inserted{de régularité locale (\unsure{de nature assez exclusive})} différente [$X_\eta$ lisse au lieu de $X_s$ d'intersection complète].}
